% Ad Familiares 16.6 --- headnote
% ad-familiares-16-06, 7 November 50 BC, Actium

\textit{Cicero to Tiro, written from Actium on the evening of
the seventh day before the Ides of November --- 7 November 50
BC --- the same day as 16.4 and 16.5, but from the far side of
the gulf, after the day's crossing from Leucas. The salutation
is the joint household one again: \textit{Tullius et Cicero et
Q. Q. Tironi s. p. d.} --- ``Tullius and Cicero and the two
Quintuses send much greeting to Tiro.''}

\textit{Cicero is candid about why he is writing for the third
time in a day: he found a courier and did not want to waste him
(\textit{magis instituti mei tenendi causa, quia nactus eram
cui darem, quam quo haberem quid scriberem}). The substance is
the same as 16.4 and 16.5 in compressed form --- love-as-care,
the long list of Tiro's services with one more to add: think
not only about your health but about your voyage. The cluster's
two great anxieties, illness and the sailing season, finally
sit together in one sentence. The closing imperative is
doubled, \textit{cura, cura te, mi Tiro}, and the formula again
is \textit{etiam atque etiam vale}.}
